\chapter{Problema y oportunidad}\label{chap:problema}

Este trabajo se desarrolla en el marco del proyecto \textit{capstone} de la maestría en ciencia de datos e inteligencia artificial de la Universidad de Ingeniería y Tecnología (UTEC), y parte de un problema que enfrentan diversas \textit{fintech} en su operación diaria. La revisión manual de documentos financieros, la validación de estados de cuenta, la identificación de datos bancarios y el cruce de información siguen siendo tareas que recaen sobre equipos de soporte, con el costo en tiempo y errores. El proyecto aborda este problema combinando ingeniería de \textit{software}, modelos de lenguaje y diseño de producto para construir una solución que extraiga y estructure automáticamente la información contenida en estos documentos, habilitando un flujo donde el propio usuario pueda completar su registro sin depender de intervención manual. En este capítulo se presenta el problema que motiva el proyecto, el contexto operativo y tecnológico en el que se inserta, la justificación de la solución propuesta y los indicadores que permitirán evaluar su impacto.


\section{Motivación y Contexto}

En muchos procesos financieros, se trabaja con documentos que presentan una alta variabilidad en su formato, contenido y calidad \cite{Zheng2023DocLevelIE, Huang2021SROIE}, lo que dificulta tratarlos de manera uniforme. Esta diversidad obliga a realizar procesos manuales de verificación, en los que se revisan datos como montos, fechas o entidades financieras, lo que incrementa los tiempos de procesamiento y genera una carga operativa constante. Además, la intervención repetida en estas tareas puede introducir inconsistencias en los registros y complica la trazabilidad de la información en escenarios de auditorías \cite{Huang2021SROIE}. La literatura comenta que los documentos visualmente ricos combinan texto y disposición espacial, lo que plantea retos específicos y ha motivado la creación de conjuntos de datos y competencias orientadas a formularios y recibos \cite{Cui2021}.

En la práctica, muchas soluciones para procesar documentos parten del uso de reconocimiento óptico de caracteres en adelante OCR, por sus siglas en inglés, que permite extraer texto desde imágenes o archivos PDF. Sin embargo, este enfoque se queda corto cuando los documentos no siguen una estructura fija, ya que no considera cómo está organizada la información dentro de la página o su \textit{layout}. Por eso, la verificación de los datos suele mantenerse como una tarea manual: el equipo de soporte revisa los documentos y valida que la información coincida con lo ingresado \cite{validacion_clabe}, lo que genera procesos más lentos y con mayor riesgo de error. Por esto, las investigaciones han explorado distintos enfoques para abordar el problema, dando lugar a estrategias que combinan tanto el contenido textual como la disposición espacial del documento. En esa línea, modelos como LayoutLM integran ambas dimensiones y han mostrado buenos resultados en tareas como formularios y recibos \cite{Xu2020LayoutLM}. De forma complementaria, se han desarrollado \textit{datasets} como DocLayNet, que buscan representar la diversidad de formatos presentes en documentos reales y mejorar la capacidad de los modelos para adaptarse a distintas estructuras \cite{Pfitzmann2022DocLayNet}. En conjunto, estos avances muestran que ya existe una base sólida para abordar documentos con estructuras variables, algo común en escenarios reales.

Al trasladar estos avances a un contexto aplicado, como el entorno mexicano donde se busca implementar la solución, el uso de inteligencia artificial en este tipo de procesos todavía es limitado. Algunas \textit{fintech} como \textit{Klar}, \textit{Konfío}, \textit{Kueski} o \textit{Clip} han incorporado IA en áreas como evaluación crediticia o detección de fraude, pero no en la extracción automática de información desde documentos bancarios institucionales \cite{klar,konfio,kueski,clip}. A nivel internacional, herramientas como \textit{Google Document AI}, \textit{Amazon Textract} o \textit{Microsoft Form Recognizer} han mostrado que el uso de \textit{deep learning} puede ser efectivo para interpretar documentos, apoyándose en avances de investigación como LayoutLM \cite{Xu2020LayoutLM} y \textit{datasets} como DocLayNet \cite{Pfitzmann2022DocLayNet}. Pese a estos avances, su adopción en contextos locales presenta limitaciones, ya sea por el costo, la dependencia de infraestructura externa o los requerimientos de cumplimiento normativo, como los establecidos en la \textit{Ley Federal de Protección de Datos Personales en Posesión de los Particulares (LFPDPPP)} \cite{lfpdppp}. Además, su desempeño puede verse afectado cuando se enfrentan a formatos específicos de cada país. En este escenario, aparece la oportunidad de desarrollar soluciones que, tomando como base estos avances, se adapten mejor a las necesidades operativas y a los documentos utilizados en la práctica.


El flujo empieza cuando la organización ofrece a instituciones educativas una plataforma digital para la gestión de pagos familiares, integrada con una pasarela de pagos como Kushki \cite{cometa_onboarding}. Esta plataforma recibe carátulas bancarias en distintos formatos, los cuales deben ser validados y estandarizados antes de su uso. Para completar este proceso, cada colegio debe proporcionar un documento oficial, como la portada de un estado de cuenta, donde se consignan datos como el nombre del titular, el número de cuenta o CLABE interbancaria\footnote{La CLABE interbancaria (Clave Bancaria Estandarizada) es un número de 18 dígitos que identifica de forma única una cuenta bancaria en México y permite realizar transferencias electrónicas seguras entre bancos, según el Banco de México. Fuente: \url{https://www.banxico.org.mx/}. En adelante, el presente documento utilizará el término CLABE como referencia al número de cuenta bancaria, por tratarse del estándar utilizado en el sistema financiero mexicano.}, el banco, entre otros. En la práctica, la validación de esta información suele implicar una revisión manual, donde el equipo de soporte verifica que los datos ingresados coincidan con el contenido del documento. Esto se vuelve especialmente complejo porque las carátulas bancarias no siguen una estructura uniforme: pueden incluir logotipos, tablas, distintas secciones y variaciones de formato según la entidad emisora. También, se suma que los documentos pueden llegar en distintos tipos de archivo, como PDF, imágenes escaneadas o fotografías tomadas con el teléfono. Como resultado, este proceso introduce fricción en la validación de cuentas bancarias receptoras de fondos y depende en gran medida de la intervención manual. En términos más amplios, este escenario representa un caso concreto de los desafíos asociados a la comprensión automática de documentos no estructurados en contextos reales. Para abordar las problemáticas antes mencionadas, el proyecto se apoya en modelos de lenguaje de gran escala como parte de una estrategia orientada a mejorar la comprensión y extracción de información en documentos. Dado que la validación bancaria requiere trabajar con datos reales y sensibles, la solución también considera mecanismos de resguardo y ofuscación durante las etapas de prueba y evaluación, con el fin de reducir la exposición innecesaria de información fuera del contexto operativo.

La validación de documentos bancarios se toma como punto de partida para desarrollar una capacidad más amplia de comprensión de documentos. A partir de este caso inicial, la solución se orienta a definir de manera explícita qué campos se desean extraer, de modo que el mismo flujo pueda aplicarse a distintos tipos de documentos que compartan una lógica similar. De esta forma, el proyecto no se queda en un único caso, sino que busca construir una base que permita adaptar el proceso de extracción y validación según el tipo de documento. Esto abre la posibilidad de trabajar con otros escenarios, como facturas, documentos de identidad o comprobantes administrativos, manteniendo el mismo enfoque pero ajustando los campos y reglas según la necesidad. Bajo esta premisa, se desarrolla un producto mínimo viable que integra las etapas principales del proceso: ingreso de documentos, definición de campos, extracción de información, validación de datos y generación de resultados estructurados. Esto permite llevar la propuesta a un entorno real, evaluar su funcionamiento en condiciones operativas y demostrar que el enfoque puede extenderse a otros casos similares. De esta manera, se establece una base clara para el desarrollo de la solución propuesta en los siguientes capítulos.

\section{Justificación}

Los principales usuarios involucrados en este proceso son, por un lado, los usuarios finales de las \textit{fintechs}, quienes deben registrar o validar sus cuentas bancarias dentro de la plataforma, y, por otro lado, los equipos de soporte u operaciones encargados de revisar, corregir o completar la información cuando se presentan inconsistencias. En un modelo B2B, el valor de la solución recae precisamente en mejorar un punto crítico del flujo operativo de la \textit{fintech}, reduciendo la fricción asociada a la carga y verificación de documentos. Cuando el proceso exige subir el archivo correcto, esperar una revisión posterior y corregir posibles errores manualmente, aumenta la probabilidad de abandono durante el \textit{onboarding} o en etapas de configuración de cuenta. En este sentido, la solución propuesta permite extraer automáticamente los campos relevantes del documento bancario y contrastarlos con la información ingresada por el usuario, facilitando una validación más ágil y consistente. Si los datos coinciden, la cuenta puede avanzar con mayor rapidez dentro del flujo; si se detectan discrepancias, estas pueden señalarse oportunamente para su corrección. De este modo, la propuesta no elimina por completo la necesidad de supervisión en todos los casos, pero sí reduce tiempos de verificación, disminuye la carga operativa sobre los equipos de soporte y contribuye a mejorar indicadores vinculados con la continuidad del proceso, como la tasa de abandono asociada a errores de documentos.

En una \textit{fintech}, el equipo de soporte se beneficia de la automatización de tareas que previamente requerían validación manual. Esto permite manejar un mayor volumen de clientes, cuentas y transacciones sin necesidad de que el equipo crezca al mismo ritmo, ayudando a mantener la calidad del servicio y a reducir la fricción en los procesos de verificación. Aunque los usuarios finales no participan de manera directa en este flujo interno, también se ven beneficiados de forma indirecta, en la medida en que esperan que las operaciones realizadas mediante la plataforma se ejecuten sobre información correctamente validada. En ese sentido, la automatización fortalece la confiabilidad general del sistema al mejorar la consistencia entre los procesos visibles para el usuario y los mecanismos internos de validación, registro y control. De esta manera, la propuesta no solo optimiza una operación interna crítica, sino que también favorece la experiencia del usuario final y se alinea con los objetivos de escalabilidad y eficiencia de la organización.


\subsection{User Stories}

De los usuarios identificados y los casos de uso descritos previamente se derivan las siguientes \textit{user stories}, las cuales recogen las necesidades y expectativas de cada actor involucrado. Gracias a esto, es posible formalizar los requisitos funcionales desde la perspectiva del usuario, asegurando que el sistema responda a las demandas de validación eficiente de cuentas bancarias, incorporación ágil de clientes y escalabilidad operativa.


\begin{itemize}
\item \textbf{US-01: Carga de documento}\\
Como usuario final, quiero subir un documento bancario de forma simple para iniciar el registro sin fricción y evitar abandonar el proceso.

\item \textbf{US-02: Extracción automática de datos}\\
Como usuario final, quiero que el sistema extraiga automáticamente los datos relevantes del documento en pocos segundos, para evitar ingresarlos manualmente y completar el registro más rápido.

\item \textbf{US-03: Revisión y corrección de datos}\\
Como usuario final, quiero revisar y corregir fácilmente los datos extraídos, para completar el registro sin necesidad de intervención del equipo de soporte.

\item \textbf{US-04: Registro de información validada}\\
Como miembro del equipo de soporte u operaciones, quiero que el sistema almacene los datos finales junto con el resultado de la extracción, para poder monitorear la calidad del proceso y detectar casos que requieran corrección manual.

\item \textbf{US-05: Reducción de intervención de soporte}\\
Como miembro del equipo de soporte u operaciones, quiero que el sistema permita a los usuarios completar y corregir sus datos de forma autónoma, para reducir la cantidad de casos que requieren intervención manual.

\item \textbf{US-06: Monitoreo de calidad de extracción}\\
Como miembro del equipo de soporte u operaciones, quiero conocer el nivel de precisión en la extracción de datos, para evaluar el desempeño del sistema y asegurar que los campos se detecten correctamente.

\item \textbf{US-07: Flujo confiable y continuo}\\
Como usuario final, quiero que el proceso de registro sea rápido y consistente, para completar el flujo sin errores ni demoras.

\end{itemize}

\subsection{Propuesta de valor y KPIs de éxito}

La propuesta de valor radica en el desarrollo de una solución basada en inteligencia artificial capaz de extraer, validar y estructurar automáticamente información contenida en distintos tipos de documentos financieros y administrativos. Esta capacidad permite atender necesidades operativas de organizaciones \textit{fintech}, especialmente en procesos asociados a \textit{compliance} y tesorería, donde la interpretación correcta y oportuna de documentos resulta crítica para la continuidad del negocio. De este modo, la solución reduce la dependencia de revisiones manuales, disminuye errores en la captura de datos y favorece la escalabilidad de procesos de documentos en contextos de alta variabilidad de formatos.

\subsection*{Objetivos SMART}
\begin{itemize}
  
    \item \textbf{Reducción del tiempo de registro:} Reducir el tiempo promedio de validación a menos de 10 segundos. Esta métrica se calculará a partir de los registros del servicio, considerando el intervalo transcurrido entre la recepción de la solicitud de procesamiento y la emisión de la respuesta correspondiente.

    \item \textbf{Extracción automatizada precisa:} Desarrollar un prototipo de API capaz de identificar correctamente los campos clave con al menos un 95\% de precisión para finales del Q2 2026 \cite{Xu2020LayoutLM}. En el caso de la CLABE, la métrica considerará como error cualquier discrepancia respecto al valor real, incluso si corresponde a un solo dígito.
  
    \item \textbf{Implementación piloto y ajuste:} Se planteó ejecutar un piloto con al menos tres organizaciones usuarias antes de marzo de 2026, con la expectativa de procesar automáticamente el 80\% de los documentos sin intervención humana. 
    
    \item \textbf{Cumplimiento regulatorio y de seguridad:} Implementar controles de seguridad y protección de datos alineados con la normativa mexicana vigente, a fin de asegurar que la solución cumpla con la Ley Federal de Protección de Datos Personales en Posesión de los Particulares \cite{lfpdppp}.
    
\end{itemize}

\subsection*{Indicadores clave (KPIs)}
\begin{itemize}
  \item Tasa de autoservicio en el registro de datos bancarios: del 75\% al 90\% de registros completados sin intervención del equipo de soporte, en un periodo de 4 a 6 meses tras la implementación.
  \item Tiempo de procesamiento por documento: $<10$ segundos por archivo procesado.
  \item Precisión de extracción por campo: $\geq$95\% en los campos principales.
  \item Casos que requieren corrección manual: $<15\%$ del total de registros procesados.
  \item Tasa de abandono durante el flujo de carga: $<10\%$ de los registros iniciados.
\end{itemize}


\noindent En primer lugar, la tasa de autoservicio en el registro de datos bancarios se medirá como la proporción de registros completados sin intervención del equipo de soporte respecto del total de registros iniciados en el flujo, con una meta de alcanzar entre el 75\% y el 90\% en un periodo de 4 a 6 meses tras la implementación. En segundo lugar, el tiempo de procesamiento por documento se calculará como el intervalo transcurrido entre la carga del archivo y la generación del resultado estructurado por parte del sistema. En tercer lugar, la precisión de extracción por campo se evaluará comparando los valores extraídos automáticamente con una verdad de terreno previamente validada para cada documento. Los casos que requieren corrección manual se medirán como el porcentaje de registros en los que un operador de soporte deba modificar o completar la información extraída. Finalmente, la tasa de abandono durante el flujo de carga se calculará como la proporción de usuarios que inician el proceso de carga de documentos pero no lo completan satisfactoriamente. Estos indicadores permitirán monitorear el desempeño del sistema y evaluar su impacto en la operación. Su cumplimiento contribuirá a validar la propuesta de valor de la solución como una herramienta capaz de reducir fricción operativa, mejorar la calidad de los datos y favorecer la escalabilidad de procesos de documentos en entornos \textit{fintech}.

